% ID: SEL-D0074-OA-E01
% Tanda S074.
\textbf{a)} Igualamos las entradas del producto con las del segundo miembro:
\[
 3(1+3x)+2y=5,\qquad 3x-y=4.
\]
La segunda ecuación da $y=3x-4$. Al sustituir en la primera se obtiene $9x+6x-8=2$, así que $\boxed{x=2/3,\ y=-2}$. En efecto, las dos entradas del producto son $9-4=5$ y $2+2=4$.

\textbf{b)} Sea $A$ la matriz dada. Para hallar su inversa, resolvemos $A(u,v,w)^t=(p,q,r)^t$. Las filas dan
\[
 u+w=p,\qquad v=q,\qquad u+2v=r.
\]
Despejando, $u=r-2q$, $v=q$ y $w=p+2q-r$. Como hay una única solución para todo segundo miembro, la matriz es invertible. Los coeficientes de $p,q,r$ en cada expresión dan las filas:
\[
 A^{-1}=\boxed{\begin{pmatrix}0&-2&1\\0&1&0\\1&2&-1\end{pmatrix}}.
\]
